%%%% CAPÍTULO 1 - INTRODUÇÃO
%%
%% Deve apresentar uma visão global da pesquisa, incluindo: breve
% histórico, importância e justificativa da escolha do tema,
%% delimitações do assunto, formulação de hipóteses e objetivos da
% pesquisa e estrutura do trabalho.

%% Título e rótulo de capítulo (rótulos não devem conter caracteres
% especiais, acentuados ou cedilha)
\chapter{Introdução}\label{cap:introducao}

A arquitetura de microsserviços consolidou-se na década de 2010 como alternativa
às aplicações monolíticas, decompondo o sistema em serviços pequenos e autônomos
que trabalham em conjunto \cite{BuildingMicroservices:2015}. Cada serviço possui
responsabilidade bem delimitada e pode ser implantado e escalado de forma
independente, o que reduz o acoplamento entre equipes e facilita a evolução
contínua do sistema \cite{MicroserviceArchitecture:2016}. Essa decomposição,
porém, torna o sistema inerentemente distribuído: chamadas antes locais passam a
atravessar a rede, trazendo desafios de latência, falhas parciais, consistência
de dados e segurança na comunicação \cite{BuildingMicroservices:2015}. Assim, a
arquitetura compensa a complexidade adicional apenas quando o sistema é grande o
bastante para se beneficiar de sua escalabilidade e autonomia.

Nesse cenário, a escolha do mecanismo de comunicação entre serviços torna-se
uma decisão arquitetural central. Entre as alternativas, destaca-se o
\textit{Remote Procedure Call} (RPC), que permite invocar procedimentos em
outro serviço como se fossem locais. O gRPC é uma das implementações mais
difundidas desse modelo: utiliza o HTTP/2 como protocolo de transporte e o
\textit{Protocol Buffers} como formato de serialização e linguagem de definição
de interfaces \cite{gRPC-Introduction:2026}. Em ambientes de produção, ele costuma ser
combinado com \textit{proxies} como o Envoy e com malhas de serviço
(\textit{service meshes}) como o Istio, que adicionam balanceamento de carga,
criptografia e observabilidade.


\section{Problema de Pesquisa}

A popularidade do gRPC faz com que ele seja frequentemente adotado como escolha
padrão, muitas vezes com base em sua reputação de eficiência. Embora existam comparações isoladas de
desempenho entre protocolos \cite{PerformanceEvaluation:2024}, é menos comum encontrar
avaliações que examinem, de forma conjunta e em um cenário realista, atributos
que frequentemente competem entre si: desempenho, disponibilidade, segurança e
tolerância a falhas. Medidas de segurança, como a criptografia mútua entre
serviços, e mecanismos de resiliência, como \textit{retries} e
\textit{circuit breakers}, por exemplo, introduzem custos que afetam o
desempenho.

Diante disso, a pergunta que orienta este trabalho é: como se comporta
um sistema de microsserviços construído com gRPC quando avaliado
simultaneamente quanto ao desempenho, à disponibilidade, à segurança e à
tolerância a falhas, e até que ponto de carga ele mantém esse comportamento?

\section{Objetivos}

O objetivo geral deste trabalho é avaliar empiricamente o gRPC em um cenário de marketplace. ara
alcançá-lo, definem-se os seguintes objetivos específicos:

\begin{itemize}
  \item projetar e implementar um sistema de microsserviços em Go, com gRPC e
    \textit{Protocol Buffers}, inspirado no \textit{system design} de um
    \textit{marketplace} de comércio eletrônico;
  \item aplicar a malha de serviço (\textit{service mesh}) Istio ao sistema, para
    prover segurança, resiliência e disponibilidade;
  \item mensurar o desempenho sob carga por meio da latência (percentis p50, p90
    e p99), da vazão (\textit{throughput}) e da saturação dos recursos de
    hardware (CPU e memória);
  \item avaliar a disponibilidade e a tolerância a falhas em cenários de falha
    induzida, por meio do tempo médio de recuperação (\textit{Mean Time To
    Recovery}, MTTR) e da taxa de erros observada durante a falha;
  \item analisar a segurança do sistema, considerando a superfície de ataque, a
    exposição a riscos e as vulnerabilidades conhecidas nas tecnologias
    empregadas;
  \item encontrar o ponto de ruptura do sistema, ou seja, o nível de carga em que
    ele deixa de atender aos critérios de desempenho e disponibilidade definidos
    para o experimento.
\end{itemize}

\section{Justificativa}

Sistemas de comércio eletrônico impõem requisitos exigentes: precisam escalar
diante de picos de demanda, permanecer disponíveis, tolerar falhas de
componentes e proteger dados sensíveis de clientes e pagamentos. Por reunir
todos esses requisitos, o \textit{marketplace} é um cenário adequado para
examinar, de forma integrada, o comportamento do gRPC em microsserviços.

Do ponto de vista prático, os resultados podem apoiar desenvolvedores,
arquitetos e engenheiros de software na decisão de adotar o gRPC, indicando
seus ganhos e custos em cada atributo de qualidade. Do ponto de vista
acadêmico, o trabalho reúne em um único experimento reproduzível dimensões que
normalmente são tratadas separadas na literatura, contribuindo para uma visão
consolidada da área.

\section{Organização do Texto}

Este trabalho está organizado em cinco capítulos, incluindo esta introdução. O
Capítulo~\ref{cap:fundamentacaoTeorica} apresenta a fundamentação teórica, com
os conceitos necessários para o entendimento do trabalho. Com base neles, o
Capítulo~\ref{cap:proposta} descreve a proposta de implementação do sistema e a
forma como ele será avaliado. Em seguida, o
Capítulo~\ref{cap:resultados} apresenta e discute os resultados dos
experimentos. Por fim, o Capítulo~\ref{cap:conclusao} reúne as conclusões.